\begin{summarybox}
	\textbf{\textcolor{CSummary}{一句话核心}}

	\textbf{硬解定理直接给出直线与椭圆联立后的二次方程系数及韦达结果，大幅减少计算量.}

	\medskip
	\textbf{\textcolor{CSummary}{统一公式}}
	\begin{itemize}[leftmargin=*, itemsep=0.5em, topsep=0.4em]
		\item 将椭圆统一写成
		      \[
			      \frac{x^2}{\alpha}+\frac{y^2}{\beta}=1
			      \qquad(\alpha>0,\ \beta>0).
		      \]
		      焦点在$x$轴时，
		      \[
			      \alpha=a^2,\qquad \beta=b^2;
		      \]
		      焦点在$y$轴时，
		      \[
			      \alpha=b^2,\qquad \beta=a^2.
		      \]

		\item 若$B\neq0$，消去$y$，得到
		      \[
			      (A^2\alpha+B^2\beta)x^2+2AC\alpha x+(C^2\alpha-B^2\alpha\beta)=0.
		      \]
		      所以
		      \[
			      x_1+x_2=-\frac{2AC\alpha}{A^2\alpha+B^2\beta},
			      \qquad
			      x_1x_2=\frac{C^2\alpha-B^2\alpha\beta}{A^2\alpha+B^2\beta}.
		      \]

		\item 若$A\neq0$，消去$x$，得到
		      \[
			      (A^2\alpha+B^2\beta)y^2+2BC\beta y+(C^2\beta-A^2\alpha\beta)=0.
		      \]
		      所以
		      \[
			      y_1+y_2=-\frac{2BC\beta}{A^2\alpha+B^2\beta},
			      \qquad
			      y_1y_2=\frac{C^2\beta-A^2\alpha\beta}{A^2\alpha+B^2\beta}.
		      \]
	\end{itemize}

	\medskip
	\textbf{\textcolor{CSummary}{使用条件}}
	\begin{itemize}[leftmargin=*, itemsep=0.5em, topsep=0.4em]
		\item \textbf{条件 1（直线一般式）：} 直线必须写为$Ax+By+C=0$形式，且$A,B$不同时为零.
		\item \textbf{条件 2（椭圆标准式）：} 椭圆必须是中心在原点、坐标轴为对称轴的标准椭圆.
		\item \textbf{条件 3（消元前提）：} 消去$y$时要求$B\neq0$；消去$x$时要求$A\neq0$.
	\end{itemize}

	\medskip
	\textbf{\textcolor{CSummary}{关键提醒}}
	\begin{itemize}[leftmargin=*, itemsep=0.5em, topsep=0.4em]
		\item \textbf{焦点轴判断：} 不要机械背$a,b$位置，先看$x^2$和$y^2$各自的分母，再确定$\alpha,\beta$.
		\item \textbf{判别式判断：} 消元后若判别式为$\Delta$，则
		      \[
			      \Delta>0
		      \]
		      表示两个不同交点；
		      \[
			      \Delta=0
		      \]
		      表示相切；
		      \[
			      \Delta<0
		      \]
		      表示没有公共点.
		\item \textbf{弦长公式前提：} 使用
		      \[
			      |AB|=\sqrt{1+k^2}\cdot\frac{\sqrt{\Delta_x}}{|p|}
		      \]
		      前，必须确认$B\neq0$，即直线斜率
		      \[
			      k=-\frac{A}{B}
		      \]
		      存在.
		\item \textbf{系数符号检查：} 代入公式前仔细核对$A,B,C$的符号，特别是$C$的符号，避免代错.
	\end{itemize}
\end{summarybox}